% Generated draft: requires mathematical and source review.
\documentclass[11pt]{article}
\usepackage[margin=25mm]{geometry}
\usepackage{fontspec}
\setmainfont{Latin Modern Roman}
\usepackage{amsmath,amssymb,mathtools,unicode-math}
\setmathfont{Latin Modern Math}
\usepackage{enumitem,needspace,xurl,hyperref}
\setlength{\parindent}{0pt}
\setlength{\parskip}{5pt}
\setlength{\emergencystretch}{3em}
\newcommand{\N}{\mathbb{N}}
\newcommand{\Z}{\mathbb{Z}}
\newcommand{\Q}{\mathbb{Q}}
\newcommand{\R}{\mathbb{R}}
\newcommand{\C}{\mathbb{C}}
\newcommand{\F}{\mathbb{F}}
\newcommand{\A}{\mathbb{A}}
\newcommand{\PP}{\mathbb{P}}
\newcommand{\E}{\mathbb{E}}
\newcommand{\eps}{\varepsilon}
\newcommand{\dd}{\,\mathrm{d}}
\begin{document}
\section*{UnsolvedMath ID: 20001934}
\textbf{Source problem code:} catalogue entry
\subsection{Definitions and mathematical statement}
Let \((S,g)\) be a closed connected Riemannian surface. In the definition reported immediately before the source question, a stationary net is a finite graph realized by geodesic edges, with each vertex of valence at least three and the balance condition
\[
 \sum_{e\ni v}\tau_{e,v}=0\in T_vS,
\]
where incident half-edges are counted and \(\tau_{e,v}\) is their outward unit tangent. The nontrivial nets used by the supplied aid have at least one genuine branching or crossing vertex; a transverse crossing of two geodesics has valence four and balances by opposite pairs. No trivalence or irreducibility condition excluding such unions is imposed in the supplied definition.

Gromov's Question 11.2.1 asks whether
\[
 \forall\text{ nonempty open }U\subset S\quad
 \exists\text{ a nontrivial stationary net }\mathcal N:
 \mathcal N\cap U\ne\varnothing.
\]
This is density of their union in the surface, not a requirement that any one finite net be dense.

Let \(\operatorname{Per}\subset SS\) denote the vectors on closed geodesic-flow orbits. The aid proves the sufficient criterion
\[
 \overline{\operatorname{Per}}=SS
 \quad\Longrightarrow\quad
 \forall U\ne\varnothing\text{ open},\ 
 \exists\gamma_1,\gamma_2\text{ distinct closed geodesic images}
\]
with a transverse intersection in \(U\). Their union, subdivided at intersections, is a stationary net whose crossing vertex in \(U\) has valence at least four. Thus the stronger statement that stationary-net crossing vertices are dense follows from periodic-direction density. The argument uses the robustness of a transverse crossing of two local geodesic segments in a surface; it does not transfer unchanged to higher dimensions.

Consequently the aid gives an affirmative answer for surfaces with Anosov geodesic flow, including every strictly negatively curved closed surface. The original background also asserts an affirmative answer for constant-curvature surfaces through density of periodic directions. These subclasses do not prove the criterion for every Riemannian metric.

A separate reported completion result is: on a closed orientable surface of positive genus, every noncontractible simple closed geodesic \(\gamma\) can be joined by a second simple closed geodesic \(\eta\) with a forced intersection. Namely one can choose an essential simple free homotopy class having positive geometric intersection number with \([\gamma]\), and a length-minimizing geodesic representative must intersect \(\gamma\). Distinct geodesic images cannot have a tangential intersection, by uniqueness of the geodesic equation. Their union is therefore a nontrivial stationary net. This completes such a given geodesic into a net, but does not ensure that such geodesics already meet every open set.

The historical source additionally says that density of the images of closed geodesics in the underlying manifold is not known in its general setting. It cites [Weinstein1970] for metrics, even bumpy metrics, on any manifold in the discussion with an open phase-space set disjoint from periodic tangent vectors. Bumpiness means nondegeneracy of closed orbits, including iterates, in the usual Poincaré-map sense. Such failure in the unit tangent bundle says nothing by itself about density of the projected closed-geodesic union in the base. The phrase ``any manifold'' has no explicit dimension qualification in the extraction. The page-number/author header is an extraction artifact, and the final announced further question is absent.
\subsection{Short English statement}
Dense periodic directions force transverse intersections of closed geodesics in every open surface region, producing stationary nets with genuine crossing vertices. This answers the question for Anosov and constant-curvature cases in the supplied account. A separate completion theorem covers each essential simple closed geodesic, but not every region of an arbitrary surface.
\subsection{Sources}
\begin{itemize}
\item Permanent UnsolvedMath ID \(20001934\), exact public code catalogue entry; local registry \texttt{aim-geometry-notes.json}, \texttt{notes[271]}, canonical tag question.
\item Keith Burns and Vladimir S. Matveev, \emph{Open Problems and Questions about Geodesics}, Question 11.2.1, attributed to Gromov, extracted source item 11.2: \url{https://aimath.org/pastworkshops/geodesicsproblems.pdf}. The supplied verification note quotes the preceding stationary-net definition. No PDF was accessed.
\item Original explicit reference [Weinstein1970], without full bibliographic details. No proposer/year metadata or original literature entry is supplied beyond the in-text attribution.
\item Research attempt 1, ``Dense periodic directions force dense stationary-net crossing vertices,'' labelled partial\_result, open at the run, proved and low novelty confidence; registry partially\_solved. The general open-status and background statements are supplied historical reports.
\end{itemize}
\end{document}
